\documentclass[12pt, a4paper]{article}
\usepackage[UTF8]{ctex}
\usepackage{geometry}
\usepackage{amsmath}
\usepackage{enumitem}
\usepackage{titlesec}
\usepackage{fancyhdr}

% 页面设置
\geometry{left=2.5cm, right=2.5cm, top=2.5cm, bottom=2.5cm}

% 标题格式
\titleformat{\section}{\large\bfseries}{\thesection}{1em}{}
\titleformat{\subsection}{\normalsize\bfseries}{\thesubsection}{1em}{}

% 页眉页脚
\pagestyle{fancy}
\fancyhf{}
\rhead{数据库原理习题集}
\lhead{专升本复习资料}
\cfoot{\thepage}


\date{}
\title{\textbf{数据库原理习题集}}
% \subtitle{专升本复习资料}
\begin{document}

\maketitle
\thispagestyle{empty}
\newpage

% ========== 第1章 绪论 ==========
\section{第1章 绪论}

\begin{enumerate}[label=\arabic*.]
	\item 以下关于数据库描述正确的是（\quad）。
	      \begin{enumerate}
		      \item 数据库是一个 DBF 文件
		      \item 数据库是一个关系
		      \item 数据库是一个结构化数据集合
		      \item 数据库是一组文件
	      \end{enumerate}
	      \textbf{答案：C}

	\item 数据独立性是指（\quad）。
	      \begin{enumerate}
		      \item 数据依赖于程序
		      \item 程序依赖于数据
		      \item 数据不依赖于程序
		      \item 程序不依赖于数据
	      \end{enumerate}
	      \textbf{答案：D}

	\item 以下关于 DBA 职责叙述中，不正确的是（\quad）。
	      \begin{enumerate}
		      \item DBA 是数据库系统超级用户，负责控制和管理各个用户访问权限
		      \item DBA 要负责监控数据库运行
		      \item DBA 要负责前端应用程序开发
		      \item DBA 要负责当数据库系统发生故障时进行恢复
	      \end{enumerate}
	      \textbf{答案：C}

	\item DBMS 是指（\quad）。
	      \begin{enumerate}
		      \item 数据库
		      \item 数据库系统
		      \item 数据库管理系统
		      \item 数据处理系统
	      \end{enumerate}
	      \textbf{答案：C}

	\item 数据冗余度低、数据共享以及较高数据独立性等特征的系统是（\quad）。
	      \begin{enumerate}
		      \item 文件系统
		      \item 数据库系统
		      \item 操作系统
		      \item 高级程序
	      \end{enumerate}
	      \textbf{答案：B}

	\item 仅次于用户和数据库之间的一层数据管理软件是（\quad）。
	      \begin{enumerate}
		      \item 数据库系统
		      \item 数据库
		      \item 管理信息系统
		      \item 数据库管理系统
	      \end{enumerate}
	      \textbf{答案：D}

	\item 数据库系统中，逻辑数据和物理数据能够相互转换，执行该功效的是（\quad）。
	      \begin{enumerate}
		      \item 操作系统
		      \item 信息管理系统
		      \item 数据库管理系统
		      \item 文件系统
	      \end{enumerate}
	      \textbf{答案：C}

	\item 数据库中对全部数据整体逻辑结构描述，称为数据库的（\quad）。
	      \begin{enumerate}
		      \item 存放模式
		      \item 子模式
		      \item 外模式
		      \item 模式
	      \end{enumerate}
	      \textbf{答案：D}

	\item 用户看到的那部分数据的局部逻辑结构描述是（\quad）。
	      \begin{enumerate}
		      \item 存放模式
		      \item 子模式（外模式）
		      \item 概念模式
		      \item 模式
	      \end{enumerate}
	      \textbf{答案：B}

	\item 文件系统和数据库系统最大的区别是（\quad）。
	      \begin{enumerate}
		      \item 数据共享
		      \item 数据独立
		      \item 数据冗余
		      \item 数据结构化
	      \end{enumerate}
	      \textbf{答案：D}

	\item 关于信息和数据，下面阐述中正确的是（\quad）。
	      \begin{enumerate}
		      \item 信息与数据，只有区分，没有联系
		      \item 信息是数据的载体
		      \item 同一信息用同一数据表示形式
		      \item 数据处理本质上就是信息处理
	      \end{enumerate}
	      \textbf{答案：D}

	\item 描述事物性质的最小数据单位是（\quad）。
	      \begin{enumerate}
		      \item 记录
		      \item 文件
		      \item 数据项
		      \item 数据库
	      \end{enumerate}
	      \textbf{答案：C}

	\item 若干记录的集合称为（\quad）。
	      \begin{enumerate}
		      \item 数据
		      \item 数据库
		      \item 数据项
		      \item 文件
	      \end{enumerate}
	      \textbf{答案：D}

	\item 数据库系统中的软件是指（\quad）。
	      \begin{enumerate}
		      \item 数据库管理系统
		      \item 应用程序
		      \item 数据库
		      \item 数据库管理员
	      \end{enumerate}
	      \textbf{答案：A}

	\item 在数据库系统组织结构中，把概念数据库与物理数据联系起来的映射是（\quad）。
	      \begin{enumerate}
		      \item 外模式/模式
		      \item 内模式/外模式
		      \item 模式/内模式
		      \item 模式/外模式
	      \end{enumerate}
	      \textbf{答案：C}

	\item 1975 年 SPARC 公布了数据库标准报告，提出了数据库（\quad）结构组织。
	      \begin{enumerate}
		      \item 一级
		      \item 二级
		      \item 三级
		      \item 四级
	      \end{enumerate}
	      \textbf{答案：C}

	\item 内模式是系统程序员用一定的（\quad）形式组织起来的一个存放文件和联系手段。
	      \begin{enumerate}
		      \item 记录
		      \item 数据
		      \item 视图
		      \item 文件
	      \end{enumerate}
	      \textbf{答案：D}

	\item 数据库系统三级结构关系，以下叙述中正确的是（\quad）。
	      \begin{enumerate}
		      \item 模式是内模式的逻辑表示
		      \item 模式是内模式的物理实现
		      \item 模式是外模式的部分抽取
		      \item 外模式是内模式的物理实现
	      \end{enumerate}
	      \textbf{答案：A}

	\item 三个模式反映了对数据库的三种不一样观点，以下说法中正确的是（\quad）。
	      \begin{enumerate}
		      \item 内模式表示了概念级数据库，表现了对数据库的总体现
		      \item 外模式表示了物理级数据库，表现了对数据库的存放观
		      \item 外模式表示了用户级数据库，表现了对数据库的用户观
		      \item 外模式表示了用户级数据库，表现了对数据库的存放观
	      \end{enumerate}
	      \textbf{答案：C}

	\item 在数据库系统组织结构中，以下（\quad）映射把用户数据库与概念数据库联系起来。
	      \begin{enumerate}
		      \item 外模式/模式
		      \item 外模式/外模式
		      \item 模式/内模式
		      \item 内模式/模式
	      \end{enumerate}
	      \textbf{答案：A}

	\item 在数据库三级模式中，只有（\quad）才是真正存放数据。
	      \begin{enumerate}
		      \item 模式
		      \item 外模式
		      \item 内模式
		      \item 用户模式
	      \end{enumerate}
	      \textbf{答案：C}

	\item 下面关于数据库管理系统阐述中，正确的是（\quad）。
	      \begin{enumerate}
		      \item 数据库管理系统是用户与应用程序的接口
		      \item 应用程序只有经过数据库管理系统才能访问数据库
		      \item 数据库管理系统用 DML 来定义三级模式
		      \item 数据库管理系统用 DDL 来实现对数据库的各种操作
	      \end{enumerate}
	      \textbf{答案：B}

	\item DBMS 经过（\quad）来定义三种模式，并将各种模式翻译成对应的目标代码。
	      \begin{enumerate}
		      \item DML
		      \item DDL
		      \item FoxPro
		      \item DBA
	      \end{enumerate}
	      \textbf{答案：B}

	\item 下面命令中，（\quad）不是 DML 基本操作。
	      \begin{enumerate}
		      \item 排序
		      \item 插入
		      \item 修改
		      \item 检索
	      \end{enumerate}
	      \textbf{答案：A}

	\item 以下关于"采取映射技术好处"的叙述中，不正确的是（\quad）。
	      \begin{enumerate}
		      \item 确保了数据独立性
		      \item 确保了数据共享
		      \item 方便了用户使用数据库
		      \item 确保了数据库开放性
	      \end{enumerate}
	      \textbf{答案：D}

	\item 数据库是指在计算机系统中按照一定数据模型组织、存放和应用的（\quad）。
	      \begin{enumerate}
		      \item 文件集合
		      \item 数据集合
		      \item 命令集合
		      \item 程序集合
	      \end{enumerate}
	      \textbf{答案：B}

	\item 数据独立性是指（\quad）。
	      \begin{enumerate}
		      \item 不会因为数据数值发生改变而影响应用程序
		      \item 不会因为系统数据存储结构和逻辑结构改变而影响程序
		      \item 不会因为程序而影响数据
		      \item 不会因为数据逻辑结构而影响数据存储结构
	      \end{enumerate}
	      \textbf{答案：B}

	\item 由计算机硬件、操作系统、DBMS、数据库、应用程序及用户等组成的一个整体称为（\quad）。
	      \begin{enumerate}
		      \item 文件系统
		      \item 数据库系统
		      \item 软件系统
		      \item 数据库管理系统
	      \end{enumerate}
	      \textbf{答案：B}

	\item 数据库最小存取单位是（\quad）。
	      \begin{enumerate}
		      \item 数据项
		      \item 字符
		      \item 记录
		      \item 文件
	      \end{enumerate}
	      \textbf{答案：C}

	\item 对于负责数据库系统、负责定义数据库内容、决定存储结构和存储策略及安全授权等工作人员是（\quad）。
	      \begin{enumerate}
		      \item 系统分析员
		      \item 数据库用户
		      \item 程序员
		      \item DBA
	      \end{enumerate}
	      \textbf{答案：D}

	\item 数据库系统支持数据独立性，这是经过（\quad）来实现。
	      \begin{enumerate}
		      \item DDL 和 DML
		      \item 定义完整性约束条件
		      \item 两级映射
		      \item DBA
	      \end{enumerate}
	      \textbf{答案：C}

	\item 以下选项不是 DBMS 功效的是（\quad）。
	      \begin{enumerate}
		      \item 数据定义
		      \item 数据操纵
		      \item 数据库运行控制
		      \item 数据编译
	      \end{enumerate}
	      \textbf{答案：D}

	\item 当数据库遭到破坏时，将其恢复到数据库破坏前的某种一致性状态，这种功效称为（\quad）。
	      \begin{enumerate}
		      \item 数据库安全性控制
		      \item 数据库完整性控制
		      \item 数据库并发控制
		      \item 数据库恢复
	      \end{enumerate}
	      \textbf{答案：D}

	\item 以下不是 DBMS 数据库运行控制功效的是（\quad）。
	      \begin{enumerate}
		      \item 并行控制
		      \item 安全性控制
		      \item 完整性控制
		      \item 故障恢复
	      \end{enumerate}
	      \textbf{答案：D}

	\item 以下关于 DBMS 叙述中，不正确的是（\quad）。
	      \begin{enumerate}
		      \item DBMS 提供数据描述语言 DDL 能够用来定义模式、外模式和内模式
		      \item 全部 DBMS 都有自含程序设计语言
		      \item 数据库运行控制包含数据库安全性、完整性、故障恢复、并发操作四个方面
		      \item DBMS 要负责实现外模式/模式以及模式/内模式之间的映射
	      \end{enumerate}
	      \textbf{答案：B}

	\item 以下关于数据库模式说法中，正确的是（\quad）。
	      \begin{enumerate}
		      \item 在应用程序中，用户使用的是内模式
		      \item 在一个数据库系统中可以有多个外模式
		      \item 模式是外模式的一个子集
		      \item 在一个数据库系统中，可以有多个内模式和外模式
	      \end{enumerate}
	      \textbf{答案：B}

	\item 定义数据库模式、数据库结构以及数据特征等功效是经过（\quad）来实现。
	      \begin{enumerate}
		      \item 数据描述语言 DDL
		      \item 数据操纵语言 DML
		      \item 程序设计语言
		      \item 机器语言
	      \end{enumerate}
	      \textbf{答案：A}

	\item 逻辑数据独立性是指（\quad）。
	      \begin{enumerate}
		      \item 模式变、用户不变
		      \item 模式变、应用程序不变
		      \item 应用程序变，模式不变
		      \item 子模式变，应用程序不变
	      \end{enumerate}
	      \textbf{答案：B}

	\item FoxPro 是一个（\quad）。
	      \begin{enumerate}
		      \item DB
		      \item DBS
		      \item DBMS
		      \item OS
	      \end{enumerate}
	      \textbf{答案：C}

	\item 三级模式间存在两种映射，它们是（\quad）。
	      \begin{enumerate}
		      \item 子模式与模式间映射，模式与内模式间映射
		      \item 子模式与内模式间映射，外模式与内模式间映射
		      \item 子模式与外模式间映射，模式与内模式间映射
		      \item 模式与内模式间映射，模式与外模式间映射
	      \end{enumerate}
	      \textbf{答案：A}

	\item FoxPro 是一个基于（\quad）的 DBMS。
	      \begin{enumerate}
		      \item 层次模型
		      \item 网状模型
		      \item 关系模式应用程序
		      \item 关系模型
	      \end{enumerate}
	      \textbf{答案：D}

	\item 模式是数据库中数据的（\quad）。
	      \begin{enumerate}
		      \item 全局物理结构
		      \item 局部物理结构
		      \item 全局逻辑结构
		      \item 局部逻辑结构
	      \end{enumerate}
	      \textbf{答案：C}

	\item 以下关于外模式叙述中，正确的是（\quad）。
	      \begin{enumerate}
		      \item 外模式是内模式的副本
		      \item 外模式是模式的子集
		      \item 外模式只能为一个应用程序所使用
		      \item 外模式与模式无关
	      \end{enumerate}
	      \textbf{答案：B}

	\item 数据冗余可能产生的问题是（\quad）。
	      \begin{enumerate}
		      \item 修改数据方便
		      \item 删除数据方便
		      \item 增加了编程复杂度
		      \item 潜在数据不一致性
	      \end{enumerate}
	      \textbf{答案：D}

	\item 以下关于 DB、DBMS、DBS 三者之间关系叙述中，正确的是（\quad）。
	      \begin{enumerate}
		      \item DB 包含 DBMS 和 DBS
		      \item DBS 包含了 DBMS 和 DB
		      \item DBMS 包含 DB 和 DBS
		      \item DB、DBMS 和 DBS 无关
	      \end{enumerate}
	      \textbf{答案：B}

	\item 下列说法不正确的是（\quad）。
	      \begin{enumerate}
		      \item 数据库减少了数据冗余
		      \item 数据库避免了一切数据重复
		      \item 数据库中的数据可以共享
		      \item 如果冗余是系统可控制的，则系统可确保更新时的一致性
	      \end{enumerate}
	      \textbf{答案：B}

	\item 以下不是数据库安全性控制方法的是（\quad）。
	      \begin{enumerate}
		      \item 设置口令
		      \item 控制用户存取权限
		      \item 数据加密
		      \item 备份
	      \end{enumerate}
	      \textbf{答案：D}

	\item 用二维表来表示实体及实体之间联系的数据模型称为（\quad）。
	      \begin{enumerate}
		      \item 实体-联系模型
		      \item 层次模型
		      \item 网状模型
		      \item 关系模型
	      \end{enumerate}
	      \textbf{答案：D}

	\item 在下述关于数据库系统叙述中，正确的是（\quad）。
	      \begin{enumerate}
		      \item 数据库中只存在数据项之间的联系
		      \item 数据库数据项之间和记录之间都存在联系
		      \item 数据库数据项之间无联系，记录之间存在联系
		      \item 数据库数据项之间和记录之间都不存在联系
	      \end{enumerate}
	      \textbf{答案：B}

	\item 数据库系统与文件系统的主要区别是（\quad）。
	      \begin{enumerate}
		      \item 数据库系统复杂，而文件系统简单
		      \item 文件系统不能处理数据冗余和数据独立性问题，而数据库系统能够处理
		      \item 文件系统只能管理程序文件，而数据库系统能够管理各种类型文件
		      \item 文件系统管理数据量较少，而数据库系统能够管理庞大数据量
	      \end{enumerate}
	      \textbf{答案：B}

	\item DBS 是采取了数据库技术的计算机系统。DBS 是一个集合体，包含数据库、计算机软硬件、应用程序和（\quad）。
	      \begin{enumerate}
		      \item 系统分析员
		      \item 程序员
		      \item 数据库管理员
		      \item 操作员
	      \end{enumerate}
	      \textbf{答案：C}

	\item 以下说法正确的是（\quad）。
	      \begin{enumerate}
		      \item 外模式、概念模式、内模式都只有一个
		      \item 模式只有一个，概念模式和内模式有多个
		      \item 外模式有多个，概念模式和内模式只有一个
		      \item 三个模式中，只有概念模式才是真正存在
	      \end{enumerate}
	      \textbf{答案：C}

	\item 要确保数据库物理数据独立性，需要修改的是（\quad）。
	      \begin{enumerate}
		      \item 模式
		      \item 模式/内模式映射
		      \item 模式/外模式映射
		      \item 内模式
	      \end{enumerate}
	      \textbf{答案：B}

	\item 以下四项中，不属于数据库系统特点的是（\quad）。
	      \begin{enumerate}
		      \item 数据共享
		      \item 数据完整性
		      \item 数据冗余较小
		      \item 数据独立性低
	      \end{enumerate}
	      \textbf{答案：D}

	\item 数据库中存放的是（\quad）。
	      \begin{enumerate}
		      \item 数据
		      \item 数据模型
		      \item 数据之间的联系
		      \item 数据以及数据之间的联系
	      \end{enumerate}
	      \textbf{答案：D}

	\item 数据库系统组织结构是（\quad）。
	      \begin{enumerate}
		      \item 三级结构三级映射
		      \item 二级结构三级映射
		      \item 三级结构二级映射
		      \item 二级结构二级映射
	      \end{enumerate}
	      \textbf{答案：C}

	\item 以下关于数据库模式说法中，正确的是（\quad）。
	      \begin{enumerate}
		      \item 三个模式中，只有外模式才是真正存在
		      \item 在应用程序中，用户使用的是外模式
		      \item 在应用程序中，用户使用的是内模式
		      \item 在应用程序中，用户使用的是概念模式
	      \end{enumerate}
	      \textbf{答案：B}

	\item 数据库（DB），数据库系统（DBS）和数据库管理系统（DBMS）三者之间的关系是（\quad）。
	      \begin{enumerate}
		      \item DBS 包括 DB 和 DBMS
		      \item DBMS 包括 DB 和 DBS
		      \item DB 包括 DBS 和 DBMS
		      \item DBS 就是 DB，也就是 DBMS
	      \end{enumerate}
	      \textbf{答案：A}

	\item 数据库系统依靠（\quad）支持了数据独立性。
	      \begin{enumerate}
		      \item 模式分级、各级之间有映像机制
		      \item 抽象数据模型，具有封装机制
		      \item 定义完整性约束条件
		      \item DDL语言和DML
	      \end{enumerate}
	      \textbf{答案：A}

	\item 数据库三级模式体系结构的划分，有利于保持数据库的（\quad）。
	      \begin{enumerate}
		      \item 数据独立性
		      \item 结构规范化
		      \item 数据安全性
		      \item 操作可行性
	      \end{enumerate}
	      \textbf{答案：A}

	\item 数据库物理存储方式的描述称为（\quad）。
	      \begin{enumerate}
		      \item 内模式
		      \item 外模式
		      \item 概念模式
		      \item 逻辑模式
	      \end{enumerate}
	      \textbf{答案：A}

	\item 在数据库的三级模式结构中，描述数据库中全体数据的全局逻辑结构和特性的是（\quad）。
	      \begin{enumerate}
		      \item 模式
		      \item 外模式
		      \item 内模式
		      \item 存储模式
	      \end{enumerate}
	      \textbf{答案：A}

	\item 数据库系统的数据独立性是指（\quad）。
	      \begin{enumerate}
		      \item 不会因为系统数据存储结构与数据逻辑结构的变化而影响应用程序
		      \item 不会因为数据的变化而影响应用程序
		      \item 不会因为存储策略的变化而影响存储结构
		      \item 不会因为某些存储结构的变化而影响其他的存储结构
	      \end{enumerate}
	      \textbf{答案：A}

	\item 关于人工管理阶段的特点，描述错误的是（\quad）。
	      \begin{enumerate}
		      \item 数据面向应用
		      \item 数据集成
		      \item 数据不保存
		      \item 应用程序管理数据
	      \end{enumerate}
	      \textbf{答案：B}

	\item 数据管理技术的发展经历了人工管理、文件系统和（\quad）三个阶段。
	      \begin{enumerate}
		      \item 数据描述阶段
		      \item 应用程序系统
		      \item 编译系统
		      \item 数据库系统
	      \end{enumerate}
	      \textbf{答案：D}

	\item 数据管理经过了手工文档、文件系统和（\quad）三个发展阶段。
	      \textbf{答案：数据库系统}

	\item 由于数据库系统在三级模式之间提供了（\quad）和（\quad）两层映像功能，这就保证了数据库系统具有较高的数据独立性。
	      \textbf{答案：外模式/模式；模式/内模式}

	\item 数据库的特点之一是数据的共享，严格地讲，这里的数据共享是指（\quad）。
	      \begin{enumerate}
		      \item 多种应用、多种语言、多个用户互相覆盖地使用数据集合
		      \item 同一个应用中的多个程序共享一个数据集合
		      \item 多个用户、同一种语言共享数据
		      \item 多个用户共享一个数据文件
	      \end{enumerate}
	      \textbf{答案：A}

	\item 在数据库的三级模式结构中内模式可以有多个。（\quad）
	      \textbf{答案：错误}

	\item 数据库系统就是数据库管理系统。（\quad）
	      \textbf{答案：错误}

	\item 模式是关系的集合。（\quad）
	      \textbf{答案：错误}

	\item 数据独立性指数据的存储与应用程序无关，数据存储结构的改变不影响应用程序的正常运行。（\quad）
	      \textbf{答案：正确}

	\item 数据是表示信息的具体形式，信息是数据表达的内容。（\quad）
	      \textbf{答案：正确}

	\item 与文件管理系统相比较，数据库系统的数据冗余度（\quad）、数据共享性（\quad）。
	      \textbf{答案：低；好}

	\item 下列关于数据库系统特点的叙述中，正确的一项是（\quad）。
	      \begin{enumerate}
		      \item 数据库系统的存储模式如有改变，概念模式无需改动
		      \item 各类用户程序均可随意地使用数据库中的各种数据
		      \item 数据库系统中概念模式改变，只需将与其有关的子模式做相应改变，不用改写应用程序
		      \item 数据一致性是指数据库中数据类型的一致
	      \end{enumerate}
	      \textbf{答案：A}

	\item 数据库系统具有的特点不包括（\quad）。
	      \begin{enumerate}
		      \item 数据冗余高
		      \item 减少应用程序开发与维护的工作量
		      \item 数据一致性
		      \item 实施统一管理与控制
	      \end{enumerate}
	      \textbf{答案：A}

	\item 数据库中不属于存储数据特点的是（\quad）。
	      \begin{enumerate}
		      \item 永久存储
		      \item 集中管理
		      \item 有组织
		      \item 可共享
	      \end{enumerate}
	      \textbf{答案：B}

	\item 实现数据在不会被相互干扰的情况下并发使用，并且在发生故障时能够对数据库进行正确恢复的是（\quad）。
	      \begin{enumerate}
		      \item 数据定义功能
		      \item 数据库的运行管理功能
		      \item 数据操纵功能
		      \item 数据库的维护功能
	      \end{enumerate}
	      \textbf{答案：B}

	\item 数据库管理系统是计算机的（\quad）。
	      \begin{enumerate}
		      \item 应用软件
		      \item 系统软件
		      \item 数据库系统
		      \item 数据库
	      \end{enumerate}
	      \textbf{答案：B}

	\item DBMS提供（\quad）来严格地定义模式。
	      \begin{enumerate}
		      \item 模式描述语言
		      \item 子模式描述语言
		      \item 内模式描述语言
		      \item 程序设计语言
	      \end{enumerate}
	      \textbf{答案：A}

	\item 数据库的核心是（\quad）。
	      \begin{enumerate}
		      \item 存储模式
		      \item 概念模式
		      \item 外部模式
		      \item 内部模式
	      \end{enumerate}
	      \textbf{答案：B}

	\item 数据库管理系统的主要目的是（\quad）。
	      \begin{enumerate}
		      \item 数据集成
		      \item 数据共享
		      \item 数据冗余小
		      \item 数据独立性高
	      \end{enumerate}
	      \textbf{答案：A}

	\item 长期储存在计算机中的有组织的、可共享的数据集合是指（\quad）。
	      \begin{enumerate}
		      \item 数据
		      \item 数据库
		      \item 数据库管理系统
		      \item 数据库系统
	      \end{enumerate}
	      \textbf{答案：B}

	\item 在计算机领域，称为数据库时代的是（\quad）。
	      \begin{enumerate}
		      \item 20世纪60年代
		      \item 20世纪70年代
		      \item 20世纪80年代
		      \item 20世纪90年代
	      \end{enumerate}
	      \textbf{答案：B}

	\item 人工管理阶段，计算机主要应用于（\quad）。
	      \begin{enumerate}
		      \item 数据集成
		      \item 过程控制
		      \item 故障恢复
		      \item 科学计算
	      \end{enumerate}
	      \textbf{答案：D}

	\item 下列不属于数据库系统三级模式结构的是（\quad）。
	      \begin{enumerate}
		      \item 模式
		      \item 外模式
		      \item 内模式
		      \item 数据模式
	      \end{enumerate}
	      \textbf{答案：D}

	\item 下列属于新一代数据库系统的是（\quad）。
	      \begin{enumerate}
		      \item 层次数据库系统
		      \item 网状数据库系统
		      \item 关系数据库系统
		      \item 面向对象数据库系统
	      \end{enumerate}
	      \textbf{答案：D}

	\item 第三代数据库系统应该是以支持（\quad）数据模型为主要特征的数据库系统。
	      \begin{enumerate}
		      \item 关系
		      \item 网状
		      \item 面向对象
		      \item 面向过程
	      \end{enumerate}
	      \textbf{答案：C}

	\item 下列属于第一代数据库系统的是（\quad）。
	      \begin{enumerate}
		      \item SYBASE
		      \item IMS
		      \item Ingres
		      \item OODBS
	      \end{enumerate}
	      \textbf{答案：B}

	\item 数据库中存储的数据的基本特点不包括（\quad）。
	      \begin{enumerate}
		      \item 永久存储
		      \item 有组织
		      \item 可共享
		      \item 易于扩展
	      \end{enumerate}
	      \textbf{答案：D}
\end{enumerate}


\end{document}
